\documentclass[11pt]{article}
\usepackage[a4paper,margin=28mm]{geometry}
\usepackage{amsmath,amssymb,amsthm,mathtools}
\usepackage{microtype}
\usepackage[hidelinks]{hyperref}

\newtheorem{theorem}{Theorem}[section]
\newtheorem{proposition}[theorem]{Proposition}
\newtheorem{corollary}[theorem]{Corollary}
\newtheorem{definition}[theorem]{Definition}
\newtheorem{remark}[theorem]{Remark}

\newcommand{\eps}{\varepsilon}
\newcommand{\D}{\mathcal D}
\newcommand{\E}{\mathcal E}
\newcommand{\Rclass}{\mathfrak R}
\newcommand{\VC}{\operatorname{VC}}
\newcommand{\PAC}{\operatorname{PAC}}
\newcommand{\blank}{\square}

\title{\(PAC_n\) and Higher-Arity VC Geometry under Fixed Finite-Monoid Distributional Typing\\
\large Typed box obstructions for CFG and MCFG incidence relations}
\author{Takayuki Kuriyama}
\date{Working draft, 2026}

\begin{document}
\maketitle

\begin{abstract}
We connect fixed finite-monoid substitutability in grammatical inference with higher-arity VC dimension and \(PAC_n\) learning. For each arity \(d\), a language induces a \((d+1)\)-ary incidence relation between a sentence context and \(d\) nonempty factors. The fixed-typing substitutability axiom forbids a labelled box pattern. Consequently, if \(h:\Sigma^*\to M\) is fixed and \(t_h=|h(\Sigma^+)|\), then the induced relation class has \(VC_2\)-dimension at most \(t_h\), while for arity \(d\ge2\) the stronger bound \(VC_{d+1}^{d-1}\le t_h\) holds. At fan-out one this bound is sharp, already for finite regular languages. The relation to \(PAC_n\) learning requires a model audit: Takeuchi's 2020 counterexample to the original full-slice \(PAC_2\) definition and the later equivalence claim of Chernikov--Towsner must first be reconciled. The unconditional results of this note are therefore stated at the higher-arity VC level, and are compared with exact positive-data reconstruction under the same observer.
\end{abstract}

\paragraph{Status.}
Exploratory project note. The typed-box upper bound and fan-out-one sharpness construction below are complete arguments. The optimal higher-arity exponent is also established below on a prime-power family of observer sizes. The exact PAC transfer is deliberately left conditional pending a definition audit. The fixed-observer exact-vs-VC separation is unconditional.

\section{Setup}

Let \(\Sigma\) be a nonempty finite alphabet, \(M\) a finite monoid, and
\[
h:\Sigma^*\to M
\]
a fixed homomorphism. Put
\[
T_h^+:=h(\Sigma^+),\qquad t_h:=|T_h^+|.
\]
Only nonempty factors are compared, as in the companion fixed-\(h\) CFG and MCFG manuscripts.

For \(d\ge1\), a sentence context is
\[
E=v_0\blank_1v_1\cdots v_{d-1}\blank_dv_d,
\]
where \(v_0,v_d\in\Sigma^*\) and \(v_1,\ldots,v_{d-1}\in\Sigma^+\).
Let \(\E_d\) be the set of these contexts. For
\(\mathbf x=(x_1,\ldots,x_d)\in(\Sigma^+)^d\), write \(E[\mathbf x]\) for
hole filling and
\[
\D_L^{(d)}(\mathbf x):=\{E\in\E_d:E[\mathbf x]\in L\},\qquad
h^{(d)}(\mathbf x):=(h(x_1),\ldots,h(x_d)).
\]

\begin{definition}[Typed tuple substitutability]
Fix \(f\ge1\). A language \(L\subseteq\Sigma^*\) is
\emph{\((f,h)\)-tuple-substitutable} if, for every \(1\le d\le f\) and
\(\mathbf x,\mathbf y\in(\Sigma^+)^d\),
\[
h^{(d)}(\mathbf x)=h^{(d)}(\mathbf y),\qquad
\D_L^{(d)}(\mathbf x)\cap\D_L^{(d)}(\mathbf y)\ne\varnothing
\]
imply
\[
\D_L^{(d)}(\mathbf x)=\D_L^{(d)}(\mathbf y).
\]
At \(d=1\) this is the ordinary two-sided fixed-\(h\) substitutability condition.
\end{definition}

Each language induces the \((d+1)\)-ary incidence relation
\[
R_L^{(d)}
:=
\{(E,x_1,\ldots,x_d)\in\E_d\times(\Sigma^+)^d:
E[x_1,\ldots,x_d]\in L\}.
\]
Let \(\Rclass_{f,h}^{(d)}\) be the family of all such relations induced by
\((f,h)\)-tuple-substitutable languages. CFG or MCFG restrictions only pass
to subclasses.

\section{The typed-box obstruction}

\begin{theorem}[Typed-box obstruction]\label{thm:typed-box}
For \(1\le d\le f\),
\[
\VC_{d+1}\!\left(\Rclass_{f,h}^{(d)}\right)
\le t_h\le |M|.
\]
The same bound holds for the subfamilies induced by fixed-\(h\) CFGs and
fixed-\(h\) fan-out-\(f\) MCFGs.
\end{theorem}

\begin{proof}
Suppose a box
\[
A_0\times A_1\times\cdots\times A_d
\subseteq \E_d\times(\Sigma^+)^d
\]
of side size \(m>t_h\) were shattered. Since \(|A_1|>t_h\), choose distinct
\(a,b\in A_1\) with \(h(a)=h(b)\). Fix arbitrary \(x_j\in A_j\) for
\(2\le j\le d\), and put
\[
\mathbf x=(a,x_2,\ldots,x_d),\qquad
\mathbf y=(b,x_2,\ldots,x_d).
\]
Then \(h^{(d)}(\mathbf x)=h^{(d)}(\mathbf y)\).

Because \(m>t_h\ge1\), choose distinct \(E_0,E_1\in A_0\). Shattering
permits a target whose trace contains
\[
(E_0,\mathbf x),\quad(E_0,\mathbf y),\quad(E_1,\mathbf x)
\]
but omits \((E_1,\mathbf y)\). Thus
\[
E_0\in\D_L^{(d)}(\mathbf x)\cap\D_L^{(d)}(\mathbf y).
\]
Typed tuple substitutability gives
\[
\D_L^{(d)}(\mathbf x)=\D_L^{(d)}(\mathbf y),
\]
contradicting the different labels at \(E_1\).
\end{proof}

\begin{remark}
The proof uses neither context-freeness nor multiple context-freeness.
The dimension bound follows from the abstract closure axiom:
equal finite type plus one common context implies equal distributions.
Grammar formalisms matter when one asks representability and effective
grammar-valued learning.
\end{remark}

\begin{theorem}[Contiguous-block strengthening]\label{thm:block-bound}
Let \(2\le d\le f\). Then
\[
  \VC_{d+1}\!\left(\Rclass_{f,h}^{(d)}\right)
  \le \left\lfloor t_h^{1/(d-1)}\right\rfloor.
\]
Equivalently, every shattered box of side size \(m\) satisfies
\[
  m^{d-1}\le t_h.
\]
\end{theorem}

\begin{proof}
Let
\[
A_0\times A_1\times\cdots\times A_d
\]
be a shattered box of side size \(m\), and fix one context
\[
E=v_0\blank_1v_1\cdots v_{d-1}\blank_dv_d\in A_0.
\]
For \((x_1,\ldots,x_{d-1})\in A_1\times\cdots\times A_{d-1}\) put
\[
z(x_1,\ldots,x_{d-1})
:=x_1v_1x_2\cdots v_{d-2}x_{d-1}\in\Sigma^+.
\]
We claim that the map
\[
(x_1,\ldots,x_{d-1})\longmapsto h(z(x_1,\ldots,x_{d-1}))
\]
is injective. Otherwise choose distinct tuples \(\mathbf x,\mathbf y\) with the
same image. Since \(m\ge2\) for any nontrivial violation of the claimed
bound, choose distinct \(a,b\in A_d\). Shattering gives a target \(L\) for which
\[
E[\mathbf x,a],\quad E[\mathbf y,a],\quad E[\mathbf x,b]\in L,
\qquad E[\mathbf y,b]\notin L.
\]
As ordinary one-factor contexts for the two block strings, use
\[
C_a=(v_0,v_{d-1}av_d),\qquad
C_b=(v_0,v_{d-1}bv_d).
\]
Then \(C_a\) belongs to both one-factor distributions of
\(z(\mathbf x)\) and \(z(\mathbf y)\), while their \(h\)-types agree. The arity-one
instance of \((f,h)\)-tuple substitutability therefore makes those distributions
equal, contradicting the two different labels at \(C_b\). Hence the map is
injective. Its domain has \(m^{d-1}\) elements and its image lies in
\(h(\Sigma^+)\), which has \(t_h\) elements.
\end{proof}

\begin{remark}
The stronger exponent is forced already by the arity-one substitutability
condition inside the higher-fan-out model. A \((d+1)\)-ary box contains, after
fixing one sentence context and one final factor coordinate, an ordinary
one-factor substitution test on a contiguous block carrying the other
\(d-1\) coordinates.
\end{remark}

\begin{definition}[Rectangular product capacity]
\label{def:rpc}
For \(r\ge1\), define \(\operatorname{rpc}_r(h)\) to be the largest \(m\)
(possibly infinite) for which there are sets
\[
 A_1,\ldots,A_r\subseteq\Sigma^+,\qquad |A_i|=m,
\]
and, when \(r\ge2\), fixed nonempty separators
\(s_1,\ldots,s_{r-1}\in\Sigma^+\), such that
\[
 (x_1,\ldots,x_r)\longmapsto
 h(x_1s_1x_2\cdots s_{r-1}x_r)
\]
is injective on \(A_1\times\cdots\times A_r\).
For \(r=1\), this is simply the maximum size of a set on which \(h\) is
injective, hence
\[
 \operatorname{rpc}_1(h)=t_h.
\]
\end{definition}

\begin{proposition}[Observer-sensitive box bound]
\label{prop:rpc-bound}
For every \(2\le d\le f\),
\[
 \VC_{d+1}\!\left(\Rclass_{f,h}^{(d)}\right)
 \le \operatorname{rpc}_{d-1}(h)
 \le \left\lfloor t_h^{1/(d-1)}\right\rfloor.
\]
Moreover, rectangular product capacity is monotone under typing refinement:
if \(\ker h'\subseteq\ker h\), then
\[
 \operatorname{rpc}_r(h')\ge\operatorname{rpc}_r(h)
\]
for every \(r\ge1\).
\end{proposition}

\begin{proof}
If a \((d+1)\)-box of side \(m\) is shattered, fix one sentence context
\[
 E=v_0\blank_1v_1\cdots v_{d-1}\blank_dv_d
\]
from its context coordinate. The proof of
Theorem~\ref{thm:block-bound} shows that
\[
 (x_1,\ldots,x_{d-1})\longmapsto
 h(x_1v_1x_2\cdots v_{d-2}x_{d-1})
\]
must be injective on the first \(d-1\) factor-coordinate sets. Therefore
\(m\le\operatorname{rpc}_{d-1}(h)\). Any injective rectangular product of
side \(m\) has \(m^{d-1}\) distinct values in \(h(\Sigma^+)\), proving the
second inequality.

For refinement, \(\ker h'\subseteq\ker h\) induces a well-defined
homomorphism from \(\operatorname{im}h'\) onto \(\operatorname{im}h\) by
\(h'(w)\mapsto h(w)\). Hence equality of two products under \(h'\) implies
equality under \(h\). Any rectangle on which the \(h\)-product map is
injective is therefore also injective for \(h'\).
\end{proof}

\begin{theorem}[Observer cardinality does not determine higher-arity VC capacity]
\label{thm:same-size-observers}
For every prime power \(q\ge4\), there are two finite observers with the
same number \(q^2\) of realized nonempty types but radically different
arity-three incidence geometry:
\[
 \VC_4\ge q
\]
for one observer, whereas
\[
 \VC_4\le1
\]
for the other.
\end{theorem}

\begin{proof}
For the first observer, take the finite-field observer
\((\mathbb F_q^2,+)\) from Theorem~\ref{thm:higher-sharp} with \(d=3\).
Its image has \(q^2\) elements and its explicit finite regular target family
shatters a \(4\)-ary box of side \(q\).

For the second, let
\[
 Z_T=\{1,z_1,\ldots,z_{T-1}\},\qquad T=q^2,
\]
where \(1\) is the identity and
\[
 z_i z_j=z_i,\qquad z_i1=z_i,\qquad 1z_j=z_j.
\]
This is a monoid. Choose an alphabet with a nonempty representative of every
element of \(Z_T\), so \(t_h=T\). We claim
\(\operatorname{rpc}_2(h)=1\). Indeed, if \(A,B\) both had at least two
elements, then \(A\) would contain a word \(x\) of nonidentity type \(z_i\).
For every separator \(s\) and every \(y\in B\),
\[
 h(xsy)=z_i,
\]
so the product map could not be injective on \(A\times B\).
Proposition~\ref{prop:rpc-bound} now gives \(\VC_4\le1\).
\end{proof}

\begin{remark}
Theorem~\ref{thm:same-size-observers} is a statistical-geometry analogue of
the project's broader \emph{observer size is not enough} theme. At higher
arity, the relevant finite resource is not only how many types the observer
has, but how its multiplication supports uniquely decodable rectangular
products.
\end{remark}

\begin{corollary}[Fan-out one]\label{cor:cfg-upper}
For fixed-\(h\) substitutable incidence relations
\[
R_L(E,x)\Longleftrightarrow E[x]\in L,
\]
one has
\[
\VC_2\le t_h.
\]
In particular, trivial typing gives \(\VC_2\le1\).
\end{corollary}

\section{Sharpness for CFG incidence relations}

\begin{proposition}[The fan-out-one bound is sharp]\label{prop:sharp-cfg}
For every \(t\ge1\), there are a finite alphabet \(\Sigma_t\), a homomorphism
\(h_t:\Sigma_t^*\to C_t\), and a family of finite regular fixed-\(h_t\)
substitutable languages whose incidence relations shatter a box of side
size \(t\). Consequently
\[
\VC_2=t=|h_t(\Sigma_t^+)|
\]
for the corresponding full fixed-\(h_t\) relation class.
\end{proposition}

\begin{proof}
Let \(C_t=\mathbb Z/t\mathbb Z\) and
\[
\Sigma_t=\{c_0,\ldots,c_{t-1},x_0,\ldots,x_{t-1}\}.
\]
Set \(h_t(c_i)=h_t(x_i)=i\bmod t\). Then
\(h_t(\Sigma_t^+)=C_t\).

For \(S\subseteq[t]\times[t]\), put
\[
L_S:=\{c_i x_j:(i,j)\in S\}.
\]
Each \(L_S\) is finite and regular. The only factors having nonempty
distribution are \(c_i\), \(x_j\), and occurring \(c_i x_j\), with
\[
\D_{L_S}(c_i)=\{(\epsilon,x_j):(i,j)\in S\},
\]
\[
\D_{L_S}(x_j)=\{(c_i,\epsilon):(i,j)\in S\},
\]
and
\[
\D_{L_S}(c_i x_j)=\{(\epsilon,\epsilon)\}.
\]
Distinct one-letter factors of the same kind have distinct types; \(c_i\)
and \(x_i\) have equal type but disjoint distributions; a one-letter factor
and a two-letter factor have disjoint distributions; and occurring
two-letter factors have identical singleton distribution. Thus the
fixed-\(h_t\) substitutability implication holds.

Let \(E_i=c_i\blank\). Then on
\[
\{E_0,\ldots,E_{t-1}\}\times\{x_0,\ldots,x_{t-1}\}
\]
we have
\[
R_{L_S}(E_i,x_j)=1\Longleftrightarrow(i,j)\in S.
\]
Arbitrary \(S\) therefore realizes every trace on the box.
\end{proof}

\section{Higher-arity sharpness at the optimal exponent}

The exponent in Theorem~\ref{thm:block-bound} is optimal. We give an
explicit rigid-skeleton construction on an infinite family of observer sizes.

\begin{theorem}[Optimal higher-arity exponent]\label{thm:higher-sharp}
Fix \(d\ge2\) and let \(q\) be a prime power with \(q\ge d+1\). There are a finite
alphabet \(\Sigma\), a homomorphism
\[
h:\Sigma^*\to (\mathbb F_q^{\,d-1},+)
\]
with
\[
t_h=q^{d-1},
\]
and a family of finite regular, \((d,h)\)-tuple-substitutable languages whose
arity-\(d\) incidence relations shatter a box of side size \(q\). Consequently
\[
\VC_{d+1}=q=t_h^{1/(d-1)}
\]
for the corresponding full relation class.
\end{theorem}

\begin{proof}
Choose vectors
\[
a_1,\ldots,a_d,b\in\mathbb F_q^{d-1}
\]
such that every \(d-1\) of them are linearly independent. Such a set exists
for \(q\ge d+1\), for example by taking \(d+1\) columns of a Vandermonde
matrix. In particular, \(a_1,\ldots,a_{d-1}\) form a basis.

Use pairwise disjoint role alphabets consisting of symbols
\[
c_i,r_i,s_{i,1},\ldots,s_{i,d-1}\quad(i\in\mathbb F_q)
\]
and
\[
x_{r,j}\quad(1\le r\le d,\ j\in\mathbb F_q).
\]
Define the letter images by
\[
h(x_{r,j})=j a_r,\qquad
h(c_i)=h(r_i)=i b,\qquad
h(s_{i,r})=0.
\]
Because \(a_1,\ldots,a_{d-1}\) span the group and a nonempty separator maps to
zero, every group element is realized by a nonempty word; hence
\(t_h=q^{d-1}\).

For
\(\mathbf i=(i_0,i_1,\ldots,i_d)\in\mathbb F_q^{d+1}\) define the rigid codeword
\[
w_{\mathbf i}
=c_{i_0}x_{1,i_1}s_{i_0,1}x_{2,i_2}\cdots
s_{i_0,d-1}x_{d,i_d}r_{i_0}.
\]
For arbitrary \(S\subseteq\mathbb F_q^{d+1}\) let
\[
L_S:=\{w_{\mathbf i}:\mathbf i\in S\}.
\]
These languages are finite and therefore regular.

We verify \((d,h)\)-tuple substitutability. Project every codeword to its
role skeleton
\[
C,X_1,S_1,X_2,\ldots,S_{d-1},X_d,R.
\]
All role alphabets are disjoint, so two tuple occurrences with a common
sentence context have the same interval positions in this skeleton. Their
common context fixes every variable \(i_r\) whose \(X_r\) position lies outside
the holes. It also fixes \(i_0\) whenever at least one of the context-specific
marker positions \(C,S_1,\ldots,S_{d-1},R\) lies outside the holes.

Consider two such tuples with equal componentwise \(h\)-types. Once \(i_0\)
is known, every component containing at most \(d-1\) variable positions is
determined by its \(h\)-type, because the corresponding subset of the
\(a_r\) is linearly independent. Thus all selected variable values are
recovered component by component unless one single interval contains all
\(d\) variable positions.

It remains only to justify recovery of \(i_0\) when all context-specific markers
lie inside the holes. With at least two holes, a nonempty gap must then
contain a variable position. The two end components contain respectively
\(c_{i_0}\) and \(r_{i_0}\). If one end component contains no variable position,
its type \(i_0b\) determines \(i_0\). Otherwise each end component contains at
most \(d-2\) variable positions; since \(b\) together with any at most \(d-2\) of
the \(a_r\) is linearly independent, equality of the end-component type again
determines \(i_0\). Once \(i_0\) is known, the preceding componentwise argument
recovers all selected variable values.

The only remaining possibility is a single interval containing all \(d\)
variable positions. It necessarily contains every internal separator
\(s_{i_0,r}\), so the factor itself fixes \(i_0\) literally; there is no variable
coordinate left outside the factor. Any two such same-type factors that
share a context therefore have the same singleton distribution. Thus in all
cases, equal tuple type plus one common context implies equal tuple
distributions, proving \((d,h)\)-tuple substitutability.

Finally, for each \(i\in\mathbb F_q\) put
\[
E_i=c_i\blank_1s_{i,1}\blank_2\cdots
s_{i,d-1}\blank_dr_i,
\]
and for coordinate \(r\) put
\[
A_r=\{x_{r,j}:j\in\mathbb F_q\}.
\]
Then
\[
R_{L_S}^{(d)}(E_{i_0},x_{1,i_1},\ldots,x_{d,i_d})=1
\quad\Longleftrightarrow\quad
(i_0,i_1,\ldots,i_d)\in S.
\]
As \(S\) is arbitrary, the \((d+1)\)-box
\(\{E_i:i\in\mathbb F_q\}\times A_1\times\cdots\times A_d\) is shattered.
The lower bound \(\VC_{d+1}\ge q\) and
Theorem~\ref{thm:block-bound} give equality.
\end{proof}

\begin{remark}
For \(d=2\) the construction can be read over the cyclic group \(C_t\) for every
\(t\ge1\): use codewords \(c_i x_j s_i y_k r_i\) with
\(h(c_i)=h(r_i)=h(x_i)=h(y_i)=i\) and \(h(s_i)=0\). Hence the \(VC_3\le t_h\)
bound is sharp for all observer sizes in arity two. The finite-field theorem
is used only to obtain a clean uniform construction in arbitrary arity.
\end{remark}

\begin{proposition}[Arbitrary observer budgets: asymptotic sharpness]
\label{prop:arbitrary-budget}
Fix \(d\ge2\), let \(T\ge1\), and put
\[
 r:=\left\lfloor T^{1/(d-1)}\right\rfloor.
\]
If \(r\ge2(d+1)\), there are a finite monoid \(M_T\) with \(|M_T|=T\), a
homomorphism \(h_T\) with
\[
 |h_T(\Sigma^+)|=T,
\]
and a family of finite regular \((d,h_T)\)-tuple-substitutable languages such
that
\[
 \VC_{d+1}>\frac r2.
\]
Together with Theorem~\ref{thm:block-bound}, the optimal higher-arity VC
dimension under an observer budget \(T\) is therefore
\[
 \Theta_d\!\left(T^{1/(d-1)}\right).
\]
\end{proposition}

\begin{proof}
Let \(q\) be the largest power of \(2\) not exceeding \(r\). Then
\(q>r/2\), and the assumption on \(r\) gives \(q\ge d+1\).
Apply Theorem~\ref{thm:higher-sharp} over \(\mathbb F_q\). Its core observer
is the additive group
\[
 G=(\mathbb F_q^{\,d-1},+)
\]
of size \(q^{d-1}\le T\), and it realizes a shattered box of side \(q\).

It remains to pad the observer without changing the target languages.
If \(k:=T-|G|>0\), adjoin elements \(z_1,\ldots,z_k\) and define
\[
 gz_i=z_i g=z_i\quad(g\in G),\qquad
 z_i z_j=z_{\max\{i,j\}}.
\]
Together with the identity of \(G\), this is a monoid \(M_T\) of size \(T\).
Add fresh letters \(e_i\) with \(h_T(e_i)=z_i\), while retaining the old
letter map into \(G\). The old alphabet realizes all elements of \(G\), and
the new letters realize all \(z_i\), so \(|h_T(\Sigma^+)|=T\).

The target languages use no \(e_i\). Hence every factor containing a fresh
letter has empty distribution and cannot create a new substitutability
premise with an occurring factor. The core languages therefore remain
\((d,h_T)\)-tuple-substitutable and retain the shattered \(q\)-box. Thus
\(\VC_{d+1}\ge q>r/2\).
\end{proof}

\section{\(PAC_n\) model audit and conditional transfer}
\label{sec:pac-audit}

The combinatorial results above do not depend on a PAC definition.  The
statistical interpretation requires more care than the preliminary version
of this note suggested.

Kuriyama--Takeuchi's 2015 definition uses the full coordinate-slice sample
\[
 D_n(\bar a)=
 \bigcup_{i<n}
 V_1\times\cdots\times V_{i-1}\times\{a_i\}\times
 V_{i+1}\times\cdots\times V_n
\]
(and unions of these sets over sampled anchors).  The 2015 paper claimed the
necessity implication
\[
 PAC_n\text{-learnable}\Longrightarrow VC_n<\infty.
\]

Takeuchi's 2020 note shows that this implication is not valid for the
original \(PAC_2\) definition as stated.  Its class
\(\mathcal C_{\mathrm{circles}}\) has infinite \(VC_2\)-dimension
(Proposition~5) but is \(PAC_2\)-learnable in the sense of the old
Definition~3 (Proposition~7); Remark~8 says that a similar example can be
made on the countable domain \(\mathbb Q^2\).  The note then restricts to
finite or countable domains and replaces the full slice by the
support-sensitive observation
\[
 D_\mu(x):=\{y\in D(x):\mu(\{y\})>0\}.
\]
For this repaired definition, it states that \(PAC_2\)-learnability again
implies finite \(VC_2\)-dimension.

Chernikov--Towsner (2025), however, formulate Definition~4.3 again with the
full coordinate slices \(D_n(\bar a)\) and state an equivalence between
finite \(VC_k\)-dimension, their packing property, and proper \(PAC_k\)
learnability.  On the face of the written definitions, this is in tension
with Takeuchi's 2020 counterexample.  There may be an additional
graded-probability-space or measurability hypothesis that resolves the
apparent conflict, but no such reconciliation is established in the present
note.

Accordingly we adopt the following conservative rule.

\begin{proposition}[Unconditional conclusion of this note]
\label{prop:unconditional-vc}
All VC bounds and sharpness results in
Theorems~\ref{thm:typed-box}, \ref{thm:block-bound},
\ref{thm:higher-sharp}, Proposition~\ref{prop:sharp-cfg}, and
Proposition~\ref{prop:arbitrary-budget} are independent of the disputed PAC
implication and remain valid under the stated substitutability definitions.
\end{proposition}

\begin{remark}[Conditional PAC transfer]
If a chosen \(PAC_k\) model is proved to satisfy
\[
 VC_k(F)<\infty\Longrightarrow F\text{ is properly }PAC_k\text{-learnable}
\]
on the countable incidence domains used here, then the VC bounds above
immediately yield the corresponding proper \(PAC_{d+1}\) learnability
statements.  Until the 2020/2025 definitions are reconciled, this is a
conditional corollary rather than a theorem of the present draft.
\end{remark}

\begin{remark}[The 2015 quantitative lower bound is quarantined]
The 2015 proof also derives a lower bound on sample complexity from a
shattered box.  Since Takeuchi (2020) explicitly identifies a failure in the
old necessity argument and repairs the observation domain, this draft does
not use that numerical lower bound until the exact proof step and model
assumptions have been re-audited.
\end{remark}

\section{Exact reconstruction versus statistical complexity}

The companion fixed-\(h\) CFG manuscript contains an operator-independent
lower bound for one fixed two-element zero-monoid observer \(h_\circ\).
Its nested finite targets \(\mathcal T_n^-\subsetneq\mathcal T_n\) have reduced
CFG representations of polynomial size and ordinary thickness \(O(n)\), while
every characteristic sample for \(\mathcal T_n\), for any set-driven operator
having characteristic samples for both targets, has encoded size at least
\[
  5\cdot 2^n-1.
\]
The typed-box theorem gives a qualitatively different complexity statement
for the incidence-relation view of the very same fixed observer.

\begin{theorem}[Fixed-observer exact-vs-VC separation]
\label{thm:observation-separation}
There is a fixed two-element monoid homomorphism \(h_\circ\) such that:

\begin{enumerate}
\item the incidence-relation class induced by
\(\mathcal C_{h_\circ}^{\mathrm{cf}}\) has uniformly bounded higher-arity
statistical geometry,
\[
  \VC_2\le2;
\]

\item no set-driven exact positive-data reconstruction operator having
characteristic samples throughout
\(\mathcal C_{h_\circ}^{\mathrm{cf}}\) admits a characteristic-data bound
polynomial in reduced CFG size and ordinary grammar thickness.  More
concretely, the companion lower-bound family forces characteristic-sample
size at least \(5\cdot2^n-1\) while both grammar size and ordinary thickness
are polynomial in \(n\).
\end{enumerate}
\end{theorem}

\begin{proof}
Part (1) is Corollary~\ref{cor:cfg-upper} with
\(|h_\circ(\Sigma^+)|\le2\).  Part (2) is exactly the set-driven
ordinary-thickness lower bound of the companion fixed-\(h\) CFG manuscript.
\end{proof}

Thus under one fixed finite observer, exact positive-data reconstruction can
require exponentially large characteristic exposure even though the
\(VC_2\)-dimension of the induced context--factor incidence class is
uniformly constant.  Any further statement converting the latter fact into
a \(PAC_2\) theorem is subject to the model audit in
Section~\ref{sec:pac-audit}.

\section{Next problems}

\begin{enumerate}
\item \textbf{Exact observer-size spectrum.} Theorem~\ref{thm:higher-sharp}
matches the root exponent on prime-power observer sizes. Determine the exact
maximum \(VC_{d+1}\) for arbitrary finite \(t_h\), not only the optimal
asymptotic exponent.

\item \textbf{Effective PAC reconstruction.} Convert the proper
relation-valued PAC learner into a computable, preferably polynomial-time,
grammar-valued learner under explicit representation assumptions.

\item \textbf{\(PAC_n\) definition audit.} Reconcile Takeuchi's 2020
counterexample with Chernikov--Towsner's 2025 Definition~4.3 and equivalence
claim before importing a statistical sample-complexity theorem.

\item \textbf{Observer refinement.} Compare the monotonic behavior of
\(VC_{d+1}\) under typing refinement with exact reconstruction cost.

\item \textbf{Repaired higher arity.} Determine the correct
support-sensitive or otherwise repaired \(PAC_n\) definition for \(n>2\),
and test whether finite \(VC_n\) is sufficient as well as necessary.
\end{enumerate}

\section{Relation to earlier grammatical-inference work}

Earlier work already gives probabilistic or stochastic learning results for
restricted CFG/MCFG families under distributional assumptions, including
Clark, Shibata--Yoshinaka, and Case--Yoshinaka--Zeugmann. The intended
contribution here is not merely probabilistic grammar learnability, but the
structural chain
\[
\text{finite-monoid typed substitutability}
\Longrightarrow
\text{forbidden higher-arity boxes}
\Longrightarrow
\text{bounded }VC_n,

\]
together with its comparison to exact reconstruction under the same observer.

\begin{thebibliography}{99}

\bibitem{KuriyamaCFG2026}
T.~Kuriyama.
\newblock Distributional learning of context-free languages under fixed finite-monoid typing.
\newblock arXiv:1409.6247v4, 2026; revised manuscript under major revision at
\emph{Theoretical Computer Science}.

\bibitem{KT2015}
T.~Kuriyama and K.~Takeuchi.
\newblock On the \(PAC_n\) learning.
\newblock RIMS K\^oky\^uroku 1938, 54--58, 2015.

\bibitem{Kob2015}
M.~Kobayashi.
\newblock A generalization of the PAC learning in product probability spaces.
\newblock RIMS K\^oky\^uroku 1938, 33--37, 2015.

\bibitem{CPT2019}
A.~Chernikov, D.~Palacin, and K.~Takeuchi.
\newblock On \(n\)-dependence.
\newblock Notre Dame Journal of Formal Logic 60(2):195--214, 2019.

\bibitem{CT2020}
A.~Chernikov and H.~Towsner.
\newblock Hypergraph regularity and higher arity VC-dimension.
\newblock arXiv:2010.00726, 2020.

\bibitem{CT2025}
A.~Chernikov and H.~Towsner.
\newblock Higher-arity PAC learning, VC dimension and packing lemma.
\newblock arXiv:2510.02420v2, 2025.

\bibitem{Takeuchi2020}
K.~Takeuchi.
\newblock On \(VC_2\)-dimension and learnability.
\newblock RIMS K\^oky\^uroku 2170:73--78, 2020.

\bibitem{ClarkEyraud2007}
A.~Clark and R.~Eyraud.
\newblock Polynomial identification in the limit of substitutable context-free languages.
\newblock Journal of Machine Learning Research 8:1725--1745, 2007.

\bibitem{Yoshinaka2011}
R.~Yoshinaka.
\newblock Efficient learning of multiple context-free languages with multidimensional substitutability from positive data.
\newblock Theoretical Computer Science 412(19):1821--1831, 2011.

\bibitem{CaseYoshinakaZeugmann2013}
J.~Case, R.~Yoshinaka, and T.~Zeugmann.
\newblock Stochastic finite learning of some mildly context-sensitive languages.
\newblock 2013.

\end{thebibliography}

\end{document}
